% پریامبل مشترک شکل‌های TikZ پایان‌نامه
% همه‌ی گیرهای xepersian + TikZ اینجا یک‌بار حل شده‌اند.
%
% ساخت هر شکل:
%   cd thesis && xelatex -output-directory=figures-tikz figures-tikz/<name>.tex
%
% قواعد نگارش داخل نودها (مهم):
%   • متن فارسی  →  \fa{...}     (bidi در جعبه‌های TikZ خودکار کار نمی‌کند)
%   • ریاضی/لاتین →  \en{...}
%   • عدد لاتین داخل ریاضی → \num{0.2}

\usepackage{fontspec}
\usepackage{amsmath,amssymb}
\usepackage{tikz}
\usetikzlibrary{arrows.meta,patterns,positioning,calc,fit,backgrounds,
                decorations.pathreplacing,shapes.geometric,shapes.misc}

\newcommand{\Sshare}{\ensuremath{S_{\mathrm{share}}}}
\newcommand{\Scrit}{\ensuremath{S_{\mathrm{crit}}}}
\newcommand{\Mf}{\ensuremath{M_f}}
\newcommand{\Mr}{\ensuremath{M_r}}

% آیکون‌ها — برداری، رنگ‌پذیر، سازگار با xepersian.
% نکته: این پکیج فقط در همین فایل‌های standalone بارگذاری می‌شود؛
% بیلد پایان‌نامه فقط PDF نهایی را \includegraphics می‌کند، پس اثری روی آن ندارد.
\usepackage{fontawesome5}

\usepackage{xepersian}

\settextfont[Scale=1.09,Extension=.ttf,Path=styles/fonts/,
  BoldFont=XB NiloofarBd,ItalicFont=XB NiloofarIt,BoldItalicFont=XB NiloofarBdIt]{XB Niloofar}
% عمدا \setdigitfont فراخوانی نمی‌شود: در شکل‌ها همه‌ی ارقام ریاضی باید لاتین بمانند.
% اگر جایی رقم فارسی لازم بود، کاراکترش را مستقیم بنویس: ۰ ۱ ۲ ۳ ۴ ۵

% ---- سه ماکرویی که همه‌ی مشکلات bidi را حل می‌کنند ----
\newcommand{\fa}[1]{\RL{#1}}                 % متن فارسی در نود TikZ
\newcommand{\en}[1]{\LR{#1}}                 % ریاضی یا لاتین در نود TikZ
\newcommand{\num}[1]{\text{\lr{#1}}}         % عدد لاتین داخل متن فارسی حالت ریاضی

% ---- پالت مشترک ----
\definecolor{colF}{RGB}{224,124,40}       % وظیفه‌ی هدف / فراموشی
\definecolor{colR}{RGB}{31,105,170}       % وظایف نگه‌داری‌شده
\definecolor{colS}{RGB}{212,92,148}       % ناحیه‌ی مشترک
\definecolor{colOK}{RGB}{34,139,84}       % حفاظت / تأیید
\definecolor{colCrit}{RGB}{106,27,154}    % زیرمجموعه‌ی بحرانی
\definecolor{colLine}{RGB}{205,205,212}
\definecolor{colHead}{RGB}{240,240,243}
\definecolor{colEmpty}{RGB}{247,247,249}
\definecolor{colInk}{RGB}{45,45,52}
\definecolor{p1}{RGB}{240,232,248}
\definecolor{p2}{RGB}{212,190,235}
\definecolor{p3}{RGB}{178,142,215}
\definecolor{p4}{RGB}{143,92,189}
\definecolor{p5}{RGB}{106,27,154}

% ---- سبک‌های مشترک ----
\tikzset{
  panel/.style   = {draw=colLine,line width=0.5pt,rounded corners=2.5pt},
  chip/.style    = {draw=#1,line width=0.7pt,rounded corners=2pt,
                    fill=#1!7,inner sep=3.5pt,align=center,font=\scriptsize},
  chip/.default  = colLine,
  flow/.style    = {-{Stealth[length=4pt,width=3pt]},line width=0.7pt,colInk},
  flowdash/.style= {-{Stealth[length=4pt,width=3pt]},line width=0.6pt,dashed,dash pattern=on 2pt off 1.6pt},
  lead/.style    = {-{Stealth[length=3pt]},line width=0.45pt,colLine},
  note/.style    = {draw=colLine,line width=0.4pt,rounded corners=2pt,
                    fill=white,inner sep=3pt,align=center,font=\scriptsize},
  stepnum/.style = {circle,draw=#1,line width=0.7pt,fill=white,inner sep=0pt,
                    minimum size=3.6mm,font=\scriptsize\bfseries,text=#1},
  stepnum/.default = colInk,
}

% سربرگ بخش: #1 عرض، #2 عنوان فارسی، #3 پایین، #4 بالا
\newcommand{\panelhead}[4]{%
  \draw[panel] (-0.35,#3) rectangle (#1,#4);
  \fill[colHead,rounded corners=2.5pt] (-0.35,{#4-0.74}) rectangle (#1,#4);
  \node[anchor=north] at ({(#1-0.35)/2},{#4-0.09}) {\small\textbf{\fa{#2}}};}

% آیکون شبکه‌ی عصبی کوچک: #1 مقیاس، #2 رنگ لبه
\newcommand{\netglyph}[2]{%
  \begin{scope}[scale=#1]
    \foreach \l/\n in {0/3,1/4,2/3}{
      \foreach \k in {1,...,\n}{
        \node[circle,draw=#2,line width=0.35pt,fill=white,inner sep=0pt,
              minimum size=1.7mm] (n\l\k) at (0.42*\l,{0.30*\k-0.15*\n}) {};}}
    \foreach \a in {1,2,3}{\foreach \b in {1,...,4}{
        \draw[#2,line width=0.15pt,opacity=0.55] (n0\a) -- (n1\b);}}
    \foreach \a in {1,...,4}{\foreach \b in {1,2,3}{
        \draw[#2,line width=0.15pt,opacity=0.55] (n1\a) -- (n2\b);}}
  \end{scope}}

% ---- سبک‌های خودکار‌اندازه (جعبه دور متن بسته می‌شود، نه برعکس) ----
\tikzset{
  % جعبه‌ی عملیاتی: عرض از محتوا می‌آید
  op/.style      = {draw=#1,line width=0.7pt,rounded corners=2.5pt,fill=white,
                    inner xsep=5pt,inner ysep=3.5pt,font=\scriptsize,align=center},
  op/.default    = colR,
  % کارت نرم: پس‌زمینه‌ی کم‌رنگ، پدینگ سخاوتمند
  card/.style 2 args = {draw=#1,line width=0.6pt,rounded corners=3.5pt,fill=#1!5,
                    inner xsep=6pt,inner ysep=5pt,align=right,font=\scriptsize,
                    text width=#2},
  cardL/.style 2 args = {draw=#1,line width=0.6pt,rounded corners=3.5pt,fill=#1!5,
                    inner xsep=6pt,inner ysep=5pt,align=center,font=\scriptsize,
                    text width=#2},
  % کارت خنثی
  plain/.style   = {draw=colLine,line width=0.6pt,rounded corners=3.5pt,fill=white,
                    inner xsep=6pt,inner ysep=5pt,align=right,font=\scriptsize,
                    text width=#1},
  plain/.default = 2.4cm,
  % نشان شماره روی لبه‌ی کارت
  badge/.style   = {circle,draw=#1,line width=0.8pt,fill=white,inner sep=0pt,
                    minimum size=4.0mm,font=\scriptsize\bfseries,text=#1},
  badge/.default = colInk,
  % عنوان بخش داخل بخش
  ttl/.style     = {font=\scriptsize\bfseries,text=#1,inner sep=1pt},
  ttl/.default   = colR,
  % زیرنویس خاکستری
  sub/.style     = {font=\scriptsize,text=colInk!62,inner sep=1pt},
}

% ---- آیکون رنگی با اندازه‌ی کنترل‌شده ----
% \ico{رنگ}{مقیاس}{\faXxx}
\newcommand{\ico}[3]{\textcolor{#1}{\scalebox{#2}{#3}}}
\usepackage{graphicx}
